% 2026-11-16 站主問：「文中的 select judges 到底是什麼情況」。
% 727、728--729、734、762、797、798 頁及註 8--9、16--17、238--240，主對話自文字層讀過（PDF 頁 + 722）。註文所引文獻未另取原件。
\chapter{法官之選任}
\label{ch:40_selection}

\begin{qaquestion}
Law\ix{persons}{law@Law, David S.} 798 頁說，任何由政府挑選法官的制度，都使政府得以「select judges that will be sympathetic to its own goals and interests」。
\begin{enumerate}
\tightlist
\item 文中談法官選任的段落有哪幾處，各在什麼脈絡？
\item 798 頁的論證與建議是什麼？
\item 這一段與導論引 Dahl 說明法院並不抗多數的段落，有何關係？
\item 人民選舉法官的建議，與本文其他部分是否相容？
\end{enumerate}
\end{qaquestion}

Law 的主張是：法院的任務是替人民監督政府，法官便不應由被監督的政府挑選；他在結論中為人民選舉法官\ix{terms}{03@\zhlidxcat{法官選任}!法官選舉}提出「a modest case」〔有限度的主張〕\pg{734 頁}。導論引 Dahl 時，同一套任命機制被用來說明法院追隨多數，與 798 頁的論證方向相反。\footcite[734]{law2009}

\section{文中各處}
\label{sec:40_selection-1}

\begin{longtable}{@{}L{0.12\linewidth}L{0.20\linewidth}L{0.23\linewidth}L{0.36\linewidth}@{}}
\caption{Law 一文論及法官選任之處：選任方式與法官別。}\\
\toprule
頁次 & 選任方式 & 法官 & 脈絡與內容 \\
\midrule
\endfirsthead
\toprule
頁次 & 選任方式 & 法官 & 脈絡與內容 \\
\midrule
\endhead
727 及註 8--9 & 人民選舉 & 美國各州法官（19 世紀紐約、伊利諾、肯塔基、維吉尼亞、俄亥俄、馬里蘭等州） & 導論。各州制憲會議代表主張民選法官，以支持積極的司法審查、強化司法獨立（引 Shugerman 博士論文） \\
728--729 及註 16--17 & 總統提名、參議院同意任命 & 美國聯邦最高法院大法官，兼及下級聯邦法院法官 & 導論。Dahl：任命程序使最高法院的政策觀點不致長期偏離立法多數；Rehnquist 認為輿論經由任命對大法官人選有發言權 \\
734 & 人民參與選任 & 未特定 & 導論預告：結論將為人民參與法官選任提出有限度的主張 \\
762 & 總統任命 & 美國聯邦最高法院大法官（假設案） & 第三部 George v. Albert。多數意見的大法官全由 George 所屬政黨的總統任命 \\
797 & 政權挑選 & 專制政權下的法院法官 & 第五部給暴君的教訓。挑選說話晦澀、冗長、好與同僚公開唱反調的法官，削弱判決作為信號的清晰度 \\
798 本文 & 政府官員任命；改由人民選舉，配合終身職或不得連任之固定任期與俸給保障 & 行使司法審查、監督政府的法院法官，未限於特定國家 & 第五部給憲法設計者的教訓。由政府任命監督者，如同狐狸任命看守雞舍者 \\
798 註 238 & 總統任命 & 美國聯邦法院法官（最高法院） & 總統偏好任命主張擴張總統權者（Moe 與 Howell） \\
798 註 239 & 人民選舉，單一長任期，政黨提名候選人（Murphy）；民選而設任期限制、不得連任（Ramseyer 與 Rasmusen） & 憲法法院法官（Murphy）；未特定（Ramseyer 與 Rasmusen 研究日本司法之專書中的一般性觀察） & 選舉與獨立可以並存的論據 \\
798 註 240 & 人民選舉，含留任選舉 & 美國州法官（明尼蘇達州；州最高法院） & 學界的疑慮：White 案 O'Connor 協同意見、Choi 等人、Shepherd \\
\bottomrule
\end{longtable}

\section{所指之法院}
\label{sec:40_selection-court}

\begin{enumerate}
\tightlist
\item
  Law 以功能界定憲法法院：負責違憲審查、終局決定憲法爭議的法院。748--749 頁並列兩種體制：歐洲與拉丁美洲多數國家由專設憲法法院\ix{terms}{02@\zhlidxcat{司法權之基礎}!憲法法院}行使司法審查，人民看該院即可；美國、加拿大採分散審查，司法體系的審級結構與上訴法院的案件篩選，使人民的注意集中於最高法院。註 89 引 Tushnet 所稱「purist Kelsenian system」〔純粹 Kelsen 式體制〕，普通法院不得就憲法問題裁判。\footcite[748--749 \& n.89]{law2009}
\item
  774 頁論法院作為焦點\ix{terms}{01@\zhlidxcat{賽局理論}!焦點}的顯著性，來自「the designation of either a specialized, high-profile constitutional court, or the nation's highest court of regular jurisdiction, to resolve constitutional disputes」〔指定專設而高知名度的憲法法院，或該國普通審判權的最高法院，解決憲法爭議〕；註 163 對照 Kelsen 模式與普通法國家的分散模式。775--776 頁另論兩個法院分享憲法解釋權的情形，註 165 舉法國（行政法院、最高法院、憲法委員會互不審查）、俄羅斯（憲法法院與最高法院）、韓國（憲法法院得宣告法律違憲而不得審查普通法院裁判，大法院得宣告行政命令違憲並廢棄下級裁判而不得宣告法律違憲）。\footcite[774--776 \& nn.163, 165]{law2009}
\item
  797 頁給憲法設計者的建議，指定「a unique and highly visible constitutional court」〔唯一而高度顯著的憲法法院〕；798 頁論選任的主詞是「courts」，以其替人民監督政府的功能為限定，未區分專設憲法法院與普通最高法院。\footcite[797--798]{law2009}
\item
  798 頁註文所引的例證分屬三種法院：註 238 為美國聯邦最高法院；註 239 Murphy 為憲法法院法官；註 240 的 White 案、Choi 等人與 Shepherd 均為美國州法院，州最高法院兼理一般民刑事上訴，憲法案件只占其案件的一部分。人民選舉法官的經驗證據來自州的普通法院，Law 的建議對象是專責違憲審查的法院，兩者的審判權範圍與案件組成不同；州法官選舉的實證能否推及專設憲法法院，註文未論。
\end{enumerate}

\section{798 頁的論證與建議}
\label{sec:40_selection-2}

\begin{enumerate}
\tightlist
\item
  前提沿用第二部的監督功能：法院若要替人民監督政府，就須不受其所監督之政府的影響。Law 寫道：「Allowing government officials to appoint the judges who will watch over them is arguably akin to allowing the foxes to appoint those who will guard the henhouse」〔讓政府官員任命監督他們的法官，可說等於讓狐狸任命看守雞舍者〕；任何由政府挑選法官的制度，都使政府得以挑選同情其目標與利益的法官，「as opposed to those of the people」〔而非同情人民的目標與利益〕。\footcite[798]{law2009}
\item
  註 238 舉例：總統可能偏好任命主張擴張總統權的法官，並引 Moe 與 Howell 的實證研究，指出最高法院傾向支持行政權的主張，在總統與國會的權限爭議中站在總統一邊。\footcite[798 n.238]{law2009}
\item
  建議：人民直接選舉法官（judicial election\ix{terms}{03@\zhlidxcat{法官選任}!法官選舉}），不把選任交給法官所要監督的人。Law 說法官選舉「may deserve more sympathetic consideration than it has received」〔或許值得比以往更同情的考量〕，並主張選舉可與其他獨立性保障並用：當選後享有終身職，或任不得連任的固定任期，並保障俸給。\footcite[798]{law2009}註 239 引 Murphy 設想憲法法院法官由人民選舉、「for a single lengthy term」〔任單一長任期〕、候選人由政黨提名；引 Ramseyer 與 Rasmusen 指出法官經選舉產生未必欠缺獨立，以任期限制使法官不得連任即可隔絕外部影響。\footcite[798 n.239]{law2009}
\item
  註 240 記錄法律學界對法官選舉的疑慮：Republican Party of Minnesota v. White\ix{cases}{02@美國!01@聯邦最高法院!2002Republican Party of Minnesota v. White（2002）@Republican Party of Minnesota v. White（2002）}, 536 U.S. 765 (2002) 中 O'Connor\ix{persons}{oconnor@O'Connor, Sandra Day} 大法官的協同意見，將該州法官公正性的問題歸咎於人民選舉法官的作法；Choi 等人的工作論文指出律師與學者的通說主張法官應由民選官員或獨立委員會任命；Shepherd 的實證研究發現，州最高法院法官的投票與決定其能否留任者的意識形態「strongly associated」〔高度相關〕。\footcite[798 n.240]{law2009}
\end{enumerate}

\section{與導論 Dahl 段的關係}
\label{sec:40_selection-3}

728--729 頁，Law 列舉政治學者對抗多數難題\ix{terms}{02@\zhlidxcat{司法權之基礎}!抗多數難題}的實證批評。Dahl\ix{persons}{dahl@Dahl, Robert A.} 的論點是，任命程序使「the policy views dominant on the Court are never for long out of line with the policy views dominant among the lawmaking majorities of the United States」〔最高法院的主流政策觀點，不會長期偏離美國立法多數的主流政策觀點〕；註 17 並引 Rehnquist\ix{persons}{rehnquist@Rehnquist, William H.}：任命大法官的權力就是引導法院政策方向的權力。\footcite[728--729 \& n.17]{law2009}

在導論中，任命使法院追隨立法多數，用以說明法院並不抗多數；在 798 頁，同一條任命的管道，是政府挑選同情自己的法官、削弱法院監督功能的途徑。兩處結論相反，Law 以委託人與代理人的區分調和：Dahl 所說的立法多數是政府，屬於代理人；法院應忠於的是人民，屬於委託人。法院追隨政府，抗多數難題固然消解，監督功能也隨之喪失，Law 因此主張由人民繞過政府直接選任法官。

\section{人民選舉法官與本文其他部分}
\label{sec:40_selection-4}

\begin{enumerate}
\tightlist
\item
  Law 在註 240 自引的 Shepherd 研究，顯示民選法官的投票與決定其留任者的意識形態高度相關。選舉法官的選民，在模型中就是委託人；法官追隨這部分選民，判決便帶有選民的黨派色彩。第 1 部第三節第 2 項所說的人民分裂，在此直接進入法院：法官由一部分選民選出，其判決對另一部分選民是否仍為焦點\ix{terms}{01@\zhlidxcat{賽局理論}!焦點}，798 頁未論。
\item
  選出政府的選民與選出法官的選民是同一批人。人民與政府在 798 頁可以區分，前提是選民在選舉政府與選舉法官時有不同的考量；註 239 Murphy 的方案由政黨提名候選人，政黨同時也是組成政府者，法官選任與政府組成又經由政黨連在一起。
\item
  786 頁的司法獨立\ix{terms}{02@\zhlidxcat{司法權之基礎}!司法獨立}，指在職法官不受政府威脅，Law 以槍手拿槍指著法官的頭為例；798 頁的選任，決定誰成為法官。第五部的賽局只把獨立列為人民的選擇\pg{787 頁}，選任方式不在賽局之內，798 頁的建議因此是賽局以外另行提出的制度主張。\footcite[786--787]{law2009}
\end{enumerate}

\section{課前可提的問題}
\label{sec:40_selection-5}

\begin{enumerate}
\tightlist
\item Law 自引 Shepherd 的實證，顯示民選法官追隨留任選民的意識形態。這與他要法院替全體人民監督政府的主張如何相容？
\item 中華民國大法官由總統提名、經立法院同意任命。依 798 頁的論證，這屬於政府挑選監督者的制度；Law 的分析能否適用於總統與立法院多數分屬不同政黨的情形？
\end{enumerate}
